\documentclass[11pt]{article}

\usepackage{amsmath,amssymb}
\usepackage{booktabs}
\usepackage{longtable}
\usepackage{array}
\usepackage[margin=1in]{geometry}
\usepackage[hidelinks]{hyperref}

\newcommand{\cA}{\mathcal{A}}
\newcommand{\barB}{\overline{B}}
\newcommand{\barBgeom}{\bar{\mathbf{B}}}
\newcommand{\barBch}{\bar{B}^{\mathrm{ch}}}
\newcommand{\Ran}{\operatorname{Ran}}
\newcommand{\MC}{\operatorname{MC}}
\newcommand{\Res}{\operatorname{Res}}
\newcommand{\ChirHoch}{\mathrm{ChirHoch}}
\DeclareMathOperator{\CoDer}{CoDer}

\title{Vol I Notation Index}
\author{}
\date{}

\begin{document}
\maketitle

This index is scoped to the notation introduced in \texttt{main.tex} and
\texttt{chapters/theory/}.
Line references record the first explicit definitional occurrence on that
surface, following Vol I document order from \texttt{main.tex}.
Earlier motivational mentions are ignored when a later line gives the first
actual symbol declaration or formula.

\bigskip

\small
\setlength{\LTleft}{0pt}
\setlength{\LTright}{0pt}
\renewcommand{\arraystretch}{1.15}

\begin{longtable}{>{\raggedright\arraybackslash}p{0.24\textwidth} >{\raggedright\arraybackslash}p{0.46\textwidth} >{\raggedright\arraybackslash}p{0.20\textwidth}}
\toprule
\textbf{Symbol} & \textbf{Meaning} & \textbf{First defined at} \\
\midrule
\endfirsthead
\toprule
\textbf{Symbol} & \textbf{Meaning} & \textbf{First defined at} \\
\midrule
\endhead
\bottomrule
\endfoot

$\barBgeom(\cA)$
& geometric bar complex
& \path{main.tex:824} \\

$\barBch(\cA)$
& abstract chiral bar complex, used when the distinction from the geometric model matters
& \path{main.tex:825} \\

$\overline{C}_n(X)$
& compactified configuration space
& \path{main.tex:826} \\

$\eta_{ij} = d\log(z_i-z_j)$
& logarithmic $1$-form used as the basic bar propagator
& \path{main.tex:827} \\

$d_0$
& genus-$0$ collision differential
& \path{main.tex:500} \\

$d_{\mathrm{fib}}$
& fiberwise curved differential on a fixed genus-$g$ curve
& \path{main.tex:501} \\

$D_g$
& total corrected differential at genus $g$
& \path{main.tex:502} \\

$\operatorname{Conv}_{\mathrm{str}}$
& strict dg Lie model of the modular convolution deformation object
& \path{main.tex:507} \\

$\mathfrak{g}_{\cA}^{\mathrm{mod}}$
& modular convolution dg Lie algebra of $\cA$
& \path{main.tex:511} \\

$\barB_X(\cA)$
& symmetric factorization bar coalgebra on unordered $\Ran(X)$; the unqualified bar notation in the geometric bar chapter
& \path{chapters/theory/bar_construction.tex:84} \\

$B^{\mathrm{ord}}_X(\cA) = T^c(s^{-1}\bar\cA)$
& ordered tensor bar; the primitive $E_1$ coalgebra whose $\Sigma_n$-coinvariants give the symmetric bar
& \path{chapters/theory/bar_construction.tex:86} \\

$\cA^i = H^*(\barB_X(\cA))$
& dual coalgebra computed by bar cohomology
& \path{chapters/theory/bar_construction.tex:90} \\

$\cA^! = ((\cA^i)^\vee)$
& Koszul dual algebra
& \path{chapters/theory/bar_construction.tex:96} \\

$\mathbb{D}_{\Ran}\,\barB_X(\cA) \simeq \cA^!_\infty$
& Verdier dual of the symmetric bar, giving the homotopy Koszul dual algebra
& \path{chapters/theory/bar_construction.tex:108} \\

$\Omega_X(\barB_X(\cA)) \xrightarrow{\sim} \cA$
& cobar-bar counit; recovers $\cA$ itself rather than $\cA^!$
& \path{chapters/theory/bar_construction.tex:135} \\

$\tau \colon \barB_X(\cA) \to \cA$
& canonical twisting morphism
& \path{chapters/theory/bar_construction.tex:142} \\

$\kappa(a,b) = a_{(1)}b$
& invariant pairing extracting the double-pole residue of the chiral product
& \path{chapters/theory/bar_construction.tex:206} \\

$\Theta_\cA = D_\cA - d_\cA^{(0)}$
& bar-intrinsic modular Maurer--Cartan element in $\mathfrak{g}_{\cA}^{\mathrm{mod}}$
& \path{chapters/theory/introduction.tex:715} \\

$\mathfrak{g}^{E_1}_\cA$
& ordered $E_1$-modular convolution algebra carrying the full $R$-matrix data before coinvariants
& \path{chapters/theory/introduction.tex:61} \\

$\Theta_\cA^{E_1} = D_\cA^{E_1} - d_\cA^{(0)}$
& ordered Maurer--Cartan lift retaining the non-averaged collision data
& \path{chapters/theory/introduction.tex:1218} \\

$\operatorname{av} \colon {\mathfrak{g}^{E_1}_\cA} \twoheadrightarrow \mathfrak{g}^{\mathrm{mod}}_\cA$
& averaging map from ordered to coinvariant modular data
& \path{chapters/theory/introduction.tex:1245} \\

$r(z) = \Res^{\mathrm{coll}}_{0,2}(\Theta_\cA^{E_1})$
& degree-$2$, genus-$0$ collision residue; the classical $r$-matrix
& \path{chapters/theory/introduction.tex:868} \\

$\kappa(\cA) = \operatorname{av}(r(z))$
& scalar modular characteristic extracted from the collision residue; the affine non-abelian Sugawara shift is handled later in the text
& \path{chapters/theory/introduction.tex:1269} \\

$\Theta_\cA^{\le r}$
& finite-order truncation of the universal Maurer--Cartan class; the levels of the shadow obstruction tower
& \path{chapters/theory/introduction.tex:1975} \\

$\Delta = 8\kappa S_4$
& critical discriminant governing finite versus infinite shadow depth
& \path{chapters/theory/introduction.tex:80} \\

$\ChirHoch^*(\cA)$
& chiral Hochschild cohomology; in the introduction it is defined as the cohomology of coderivations of the bar coalgebra
& \path{chapters/theory/introduction.tex:694} \\

$Z^{\mathrm{der}}_{\mathrm{ch}}(\cA) = \operatorname{HH}^*(\cA,\cA)$
& derived chiral center; the bulk object computed by Hochschild cochains
& \path{main.tex:787} \\

$\omega_g = c_1(\mathbb{E})$
& Hodge class governing fiberwise curvature
& \path{chapters/theory/introduction.tex:53} \\

$\lambda_g = c_g(\mathbb{E})$
& top Chern, or Mumford, class of the Hodge bundle
& \path{chapters/theory/higher_genus_foundations.tex:188} \\

$\mathrm{obs}_g(\cA) = \kappa(\cA)\cdot\lambda_g$
& scalar genus-$g$ obstruction class on the uniform-weight lane
& \path{chapters/theory/higher_genus_foundations.tex:186} \\

$\Lambda = \sum_{g \ge 1} \lambda_g$
& universal Hodge-class series used in the scalar shadow package
& \path{chapters/theory/higher_genus_modular_koszul.tex:4670} \\

$\mathcal{C}_{\cA} = (\Theta_\cA,\kappa(\cA),\Delta_\cA,\Pi_\cA,\mathcal{H}_\cA)$
& full modular characteristic package: universal MC class, scalar characteristic, discriminant, periodicity profile, and Hilbert series package
& \path{chapters/theory/higher_genus_modular_koszul.tex:2652} \\

$\kappa_{\mathrm{ch}},\kappa_{\mathrm{cat}},\kappa_{\mathrm{BKM}},\kappa_{\mathrm{fiber}}$
& Vol III subscripted modular characteristic; bare $\kappa$ is forbidden on that volume. $\kappa_{\mathrm{ch}}$ records the chiral shadow of the CY-to-chiral functor, $\kappa_{\mathrm{cat}}$ the categorified Euler trace, $\kappa_{\mathrm{BKM}}$ the Borcherds weight $c_N(0)/2$, and $\kappa_{\mathrm{fiber}}$ the lattice/fiber correction
& \path{Vol III landscape_census_cy.tex} \\

$Y_\hbar(\fg)$, $Y^{\mathrm{dg}}_\hbar(\fg)$, $Y(\fg)^{\mathrm{ch}}$, $Y^{\mathrm{spectral}}(\fg)$
& four Yangian types: classical $E_1$-topological on $\mathbb{R}$; dg-shifted on the formal disk; $E_1$-chiral as $\mathcal{D}$-module on a curve; spectral factorization algebra on $\mathbb{A}^1_u$. Conflation is a type error
& \path{chapters/theory/e1_primacy_ordered_bar.tex} \\

$\mathrm{HH}^*_{\mathrm{top}},\ \ChirHoch^*,\ \mathrm{HH}^*_{\mathrm{cat}}$
& the three Hochschild theories: topological (Deligne, $E_1\to E_2$), chiral (Theorem~H, $E_\infty$-chiral input, concentrated in degrees $\{0,1,2\}$), and categorical (dg category input, $E_2$ with CY shifted Poisson). The geometry determines which
& \path{appendices/concordance.tex} \\

$\mathrm{SC}^{\mathrm{ch,top}}$
& Swiss-cheese chiral/topological operad controlling the derived chiral centre pair $(\mathrm{C}^{\bullet}_{\mathrm{ch}}(\cA,\cA),\cA)$. Two coloured, not self-dual; Dunn additivity does not apply
& \path{chapters/theory/en_chiral_operadic_circle.tex} \\

$\barB^{\mathrm{ord}}(\cA)\ \text{vs.}\ \barB^{\Sigma}(\cA)$
& ordered bar (primitive $E_1$-coalgebra with deconcatenation, carrying the $R$-matrix and Yangian) versus symmetric bar ($\Sigma_n$-coinvariant shadow carrying only $\kappa$). The ordered bar generates; the symmetric bar averages
& \path{chapters/theory/e1_primacy_ordered_bar.tex} \\

$\mathrm{av}\colon \mathfrak{g}^{E_1}_\cA \twoheadrightarrow \mathfrak{g}^{\mathrm{mod}}_\cA$
& lossy $\Sigma_n$-coinvariant projection from ordered to coinvariant modular data. At degree~$2$ on non-abelian affine Kac--Moody it recovers $\kappa_{\mathrm{dp}}$; the full $\kappa$ adds the Sugawara shift $\dim(\fg)/2$
& \path{chapters/theory/introduction.tex:1245} \\

$C_\cA \;=\; 6$
& universal coefficient-growth asymptotic on non-logarithmic class~M with Virasoro subalgebra; inverse radius of convergence of $\Sigma_\cA(z) = \sum_{r\ge 0} A_r z^r$. Equivalent data: the closed form $\Sigma_{\Vir}(z) = 4z^3/9 - z^2/9 + z/27 - \log(1+6z)/162$
& \path{chapters/theory/shadow_tower_higher_coefficients.tex} \\

$\Delta = 8\kappa S_4$
& critical discriminant; $\Delta = 0$ iff the shadow tower is of finite depth
& \path{chapters/theory/introduction.tex:80} \\

G/L/C/M
& the four depth classes of the shadow quadrichotomy: Gaussian $r=2$ (Heisenberg), Lie $r=3$ (affine Kac--Moody), contact $r=4$ ($\beta\gamma$), mixed $r=\infty$ (Virasoro, principal $\cW_N$)
& \path{chapters/theory/classification.tex} \\

\end{longtable}

\noindent\emph{Cross-volume note.} Vol~II and Vol~III add families of
notation that do not appear above: Vol~II uses the lambda-bracket
convention (so $S_2^{\Vol\,\text{II}}$ equals $c/12$, not $c/2$; the
bridge is $S_2^{\Vol\,\text{I}} = 6\cdot S_2^{\Vol\,\text{II}}$, see
\texttt{wave\_supervisory\_q\_convention\_bridge.md}); Vol~III deploys
the CY-to-chiral functor $\Phi$ and the universal trace identity
relating $\kappa_{\mathrm{BKM}} = c_N(0)/2$ to $\kappa_{\mathrm{ch}}$.
The bare $\kappa$ convention of Vol~I is replaced there by the
subscripted registry \{$\mathrm{ch},\mathrm{cat},\mathrm{BKM},
\mathrm{fiber}\}$ (approved, closed set).

\end{document}
